\chapter{PREFACE}

Non-Euclidean Geometry is now recognized as an important branch of
Mathematics. Those who teach Geometry should have some knowledge of
this subject, and all who are interested in Mathematics will find
much to stimulate them and much for them to enjoy in the novel
results and views that it presents.

This book is an attempt to give a simple and direct account of the
Non-Euclidean Geometry, and one which presupposes but little
knowledge of Mathematics. The first three chapters assume a
knowledge of only Plane and Solid Geometry and Trigonometry, and the
entire book can be read by one who has taken the mathematical
courses commonly given in our colleges.

No special claim to originality can be made for what is published
here. The propositions have long been established, and in various
ways. Some of the proofs may be new, but others, as already given by
writers on this subject, could not be improved. These have come to
me chiefly through the translations of Professor George Bruce
Halsted of the University of Texas.

I am particularly indebted to my friend, Arnold B.~Chace, Sc.D., of
Valley Falls, R.~I., with whom I have studied and discussed the
subject.\medskip

\hfill HENRY P.~MANNING.

{\footnotesize \textsc{Providence}, January, 1901.}

%% CONTENTS
%% INTRODUCTION
%% CHAPTER I: PANGEOMETRY
%%    I. PROPOSITIONS DEPENDING ONLY ON THE PRINCIPLE OF
%%       SUPERPOSITION
%%   II. PROPOSITIONS WHICH ARE TRUE FOR RESTRICTED FIGURES
%%  III. THE THREE HYPOTHESES
%%
%% CHAPTER II: THE HYPERBOLIC GEOMETRY
%%    I. PARALLEL LINES
%%   II. BOUNDARY-CURVES AND SURFACES, AND EQUIDISTANT-CURVES
%%       AND SURFACES
%%  III. TRIGONOMETRICAL FORMULÆ
%%
%% CHAPTER III: THE ELLIPTIC GEOMETRY
%%
%% CHAPTER IV: ANALYTIC NON-EUCLIDEAN GEOMETRY
%%    I. HYPERBOLIC ANALYTIC GEOMETRY
%%   II. ELLIPTIC ANALYTIC GEOMETRY
%%  III. ELLIPTIC SOLID ANALYTIC GEOMETRY
%%
%% HISTORICAL NOTE
\tableofcontents
\mainmatter
